\subsection{\label{sec.details.gf.lim}Limits of FPSs}

\begin{fineprint}
Let us now prove a few facts stated in Section \ref{sec.gf.lim}.

\begin{proof}
[Detailed proof of Lemma \ref{lem.fps.lim.xn-equiv}.]We have $\lim
\limits_{i\rightarrow\infty}f_{i}=f$. In other words, the sequence $\left(
f_{i}\right)  _{i\in\mathbb{N}}$ coefficientwise stabilizes to $f$. In other
words, for each $n\in\mathbb{N}$,%
\begin{equation}
\text{the sequence }\left(  \left[  x^{n}\right]  f_{i}\right)  _{i\in
\mathbb{N}}\text{ stabilizes to }\left[  x^{n}\right]  f
\label{pf.lem.fps.lim.xn-equiv.stab}%
\end{equation}
(by the definition of \textquotedblleft coefficientwise
stabilizing\textquotedblright).

Now, let $n\in\mathbb{N}$. Furthermore, let $k\in\left\{  0,1,\ldots
,n\right\}  $. Then, the sequence $\left(  \left[  x^{k}\right]  f_{i}\right)
_{i\in\mathbb{N}}$ stabilizes to $\left[  x^{k}\right]  f$ (by
(\ref{pf.lem.fps.lim.xn-equiv.stab}), applied to $k$ instead of $n$). In other
words, there exists some $N\in\mathbb{N}$ such that%
\[
\text{all integers }i\geq N\text{ satisfy }\left[  x^{k}\right]  f_{i}=\left[
x^{k}\right]  f
\]
(by the definition of \textquotedblleft stabilizing\textquotedblright). Let us
denote this $N$ by $N_{k}$. Thus,
\begin{equation}
\text{all integers }i\geq N_{k}\text{ satisfy }\left[  x^{k}\right]
f_{i}=\left[  x^{k}\right]  f. \label{pf.lem.fps.lim.xn-equiv.Nk}%
\end{equation}


Forget that we fixed $k$. Thus, for each $k\in\left\{  0,1,\ldots,n\right\}
$, we have defined an integer $N_{k}\in\mathbb{N}$ for which
(\ref{pf.lem.fps.lim.xn-equiv.Nk}) holds. Altogether, we have thus defined
$n+1$ integers $N_{0},N_{1},\ldots,N_{n}\in\mathbb{N}$.

Let us set $P:=\max\left\{  N_{0},N_{1},\ldots,N_{n}\right\}  $. Then, of
course, $P\in\mathbb{N}$.

Now, let $i\geq P$ be an integer. Then, for each $k\in\left\{  0,1,\ldots
,n\right\}  $, we have $i\geq P=\max\left\{  N_{0},N_{1},\ldots,N_{n}\right\}
\geq N_{k}$ (since $k\in\left\{  0,1,\ldots,n\right\}  $) and therefore
$\left[  x^{k}\right]  f_{i}=\left[  x^{k}\right]  f$ (by
(\ref{pf.lem.fps.lim.xn-equiv.Nk})). In other words, each $k\in\left\{
0,1,\ldots,n\right\}  $ satisfies $\left[  x^{k}\right]  f_{i}=\left[
x^{k}\right]  f$. Renaming the variable $k$ as $m$ in this result, we obtain
the following:%
\[
\text{Each }m\in\left\{  0,1,\ldots,n\right\}  \text{ satisfies }\left[
x^{m}\right]  f_{i}=\left[  x^{m}\right]  f\text{.}%
\]
In other words, $f_{i}\overset{x^{n}}{\equiv}f$ (by the definition of $x^{n}%
$-equivalence\footnote{i.e., Definition \ref{def.fps.xneq}}).

Forget that we fixed $i$. We thus have shown that all integers $i\geq P$
satisfy $f_{i}\overset{x^{n}}{\equiv}f$. Hence, there exists some integer
$N\in\mathbb{N}$ such that
\[
\text{all integers }i\geq N\text{ satisfy }f_{i}\overset{x^{n}}{\equiv}f
\]
(namely, $N=P$). This proves Lemma \ref{lem.fps.lim.xn-equiv}.
\end{proof}

\begin{proof}
[Detailed proof of Proposition \ref{prop.fps.lim.sum-prod}.]Recall that
$\lim\limits_{i\rightarrow\infty}f_{i}=f$. Hence, Lemma
\ref{lem.fps.lim.xn-equiv} shows that there exists some integer $N\in
\mathbb{N}$ such that
\[
\text{all integers }i\geq N\text{ satisfy }f_{i}\overset{x^{n}}{\equiv}f.
\]
Let us denote this $N$ by $K$. Hence, all integers $i\geq K$ satisfy
$f_{i}\overset{x^{n}}{\equiv}f$.

Thus, we have found an integer $K\in\mathbb{N}$ such that%
\begin{equation}
\text{all integers }i\geq K\text{ satisfy }f_{i}\overset{x^{n}}{\equiv}f.
\label{prop.fps.lim.sum-prod.K}%
\end{equation}
Similarly, using $\lim\limits_{i\rightarrow\infty}g_{i}=g$, we can find an
integer $L\in\mathbb{N}$ such that
\begin{equation}
\text{all integers }i\geq L\text{ satisfy }g_{i}\overset{x^{n}}{\equiv}g.
\label{prop.fps.lim.sum-prod.L}%
\end{equation}
Consider these $K$ and $L$. Let us furthermore set $P:=\max\left\{
K,L\right\}  $. Thus, $P\in\mathbb{N}$.

We shall now show that each integer $i\geq P$ satisfies $\left[  x^{n}\right]
\left(  f_{i}g_{i}\right)  =\left[  x^{n}\right]  \left(  fg\right)  $.

Indeed, let $i\geq P$ be an integer. Then, (\ref{prop.fps.lim.sum-prod.K})
yields $f_{i}\overset{x^{n}}{\equiv}f$ (since $i\geq P=\max\left\{
K,L\right\}  \geq K$), whereas (\ref{prop.fps.lim.sum-prod.L}) yields
$g_{i}\overset{x^{n}}{\equiv}g$ (since $i\geq P=\max\left\{  K,L\right\}  \geq
L$). Hence, we obtain%
\[
f_{i}g_{i}\overset{x^{n}}{\equiv}fg
\]
(by (\ref{eq.thm.fps.xneq.props.b.*}), applied to $a=f_{i}$, $b=f$, $c=g_{i}$
and $d=g$). In other words,
\[
\text{each }m\in\left\{  0,1,\ldots,n\right\}  \text{ satisfies }\left[
x^{m}\right]  \left(  f_{i}g_{i}\right)  =\left[  x^{m}\right]  \left(
fg\right)
\]
(by the definition of $x^{n}$-equivalence). Applying this to $m=n$, we find%
\[
\left[  x^{n}\right]  \left(  f_{i}g_{i}\right)  =\left[  x^{n}\right]
\left(  fg\right)  .
\]


Forget that we fixed $i$. We thus have shown that all integers $i\geq P$
satisfy $\left[  x^{n}\right]  \left(  f_{i}g_{i}\right)  =\left[
x^{n}\right]  \left(  fg\right)  $. Hence, there exists some $N\in\mathbb{N}$
such that%
\[
\text{all integers }i\geq N\text{ satisfy }\left[  x^{n}\right]  \left(
f_{i}g_{i}\right)  =\left[  x^{n}\right]  \left(  fg\right)
\]
(namely, $N=P$). In other words,%
\[
\text{the sequence }\left(  \left[  x^{n}\right]  \left(  f_{i}g_{i}\right)
\right)  _{i\in\mathbb{N}}\text{ stabilizes to }\left[  x^{n}\right]  \left(
fg\right)
\]
(by the definition of \textquotedblleft stabilizes\textquotedblright).

Forget that we fixed $n$. We thus have shown that for each $n\in\mathbb{N}$,
the sequence $\left(  \left[  x^{n}\right]  \left(  f_{i}g_{i}\right)
\right)  _{i\in\mathbb{N}}$ stabilizes to $\left[  x^{n}\right]  \left(
fg\right)  $. In other words, the sequence $\left(  f_{i}g_{i}\right)
_{i\in\mathbb{N}}$ coefficientwise stabilizes to $fg$ (by the definition of
\textquotedblleft coefficientwise stabilizing\textquotedblright). In other
words, $\lim\limits_{i\rightarrow\infty}\left(  f_{i}g_{i}\right)  =fg$.

The same argument -- but using (\ref{eq.thm.fps.xneq.props.b.+}) instead of
(\ref{eq.thm.fps.xneq.props.b.*}) -- shows that $\lim\limits_{i\rightarrow
\infty}\left(  f_{i}+g_{i}\right)  =f+g$. Thus, the proof of Proposition
\ref{prop.fps.lim.sum-prod} is complete.
\end{proof}

\begin{proof}
[Detailed proof of Proposition \ref{prop.fps.lim.sum-quot}.]First, let us show
that $g$ is invertible:

\begin{statement}
\textit{Claim 1:} The FPS $g\in K\left[  \left[  x\right]  \right]  $ is invertible.
\end{statement}

\begin{proof}
[Proof of Claim 1.]We have assumed that $\lim\limits_{i\rightarrow\infty}%
g_{i}=g$. In other words, the sequence $\left(  g_{i}\right)  _{i\in
\mathbb{N}}$ coefficientwise stabilizes to $g$. In other words, for each
$n\in\mathbb{N}$, the sequence $\left(  \left[  x^{n}\right]  g_{i}\right)
_{i\in\mathbb{N}}$ stabilizes to $\left[  x^{n}\right]  g$ (by the definition
of \textquotedblleft coefficientwise stabilizing\textquotedblright). Applying
this to $n=0$, we see that the sequence $\left(  \left[  x^{0}\right]
g_{i}\right)  _{i\in\mathbb{N}}$ stabilizes to $\left[  x^{0}\right]  g$. In
other words, there exists some $N\in\mathbb{N}$ such that%
\[
\text{all integers }i\geq N\text{ satisfy }\left[  x^{0}\right]  g_{i}=\left[
x^{0}\right]  g
\]
(by the definition of \textquotedblleft stabilizes\textquotedblright).
Consider this $N$. Thus, all integers $i\geq N$ satisfy $\left[  x^{0}\right]
g_{i}=\left[  x^{0}\right]  g$. Applying this to $i=N$, we obtain $\left[
x^{0}\right]  g_{N}=\left[  x^{0}\right]  g$.

However, each FPS $g_{i}$ is invertible (by assumption). Hence, in particular,
the FPS $g_{N}$ is invertible. By Proposition \ref{prop.fps.invertible}
(applied to $a=g_{N}$), this entails that its constant term $\left[
x^{0}\right]  g_{N}$ is invertible in $K$. In other words, $\left[
x^{0}\right]  g$ is invertible in $K$ (since $\left[  x^{0}\right]
g_{N}=\left[  x^{0}\right]  g$).

But the FPS $g$ is invertible if and only if its constant term $\left[
x^{0}\right]  g$ is invertible in $K$ (by Proposition
\ref{prop.fps.invertible}, applied to $a=g$). Hence, $g$ is invertible (since
its constant term $\left[  x^{0}\right]  g$ is invertible in $K$). This proves
Claim 1.
\end{proof}

It remains to prove that $\lim\limits_{i\rightarrow\infty}\dfrac{f_{i}}{g_{i}%
}=\dfrac{f}{g}$.

Recall that $\lim\limits_{i\rightarrow\infty}f_{i}=f$. Hence, Lemma
\ref{lem.fps.lim.xn-equiv} shows that there exists some integer $N\in
\mathbb{N}$ such that
\[
\text{all integers }i\geq N\text{ satisfy }f_{i}\overset{x^{n}}{\equiv}f.
\]
Let us denote this $N$ by $K$. Hence, all integers $i\geq K$ satisfy
$f_{i}\overset{x^{n}}{\equiv}f$.

Thus, we have found an integer $K\in\mathbb{N}$ such that%
\begin{equation}
\text{all integers }i\geq K\text{ satisfy }f_{i}\overset{x^{n}}{\equiv}f.
\label{pf.prop.fps.lim.sum-quot.K}%
\end{equation}
Similarly, using $\lim\limits_{i\rightarrow\infty}g_{i}=g$, we can find an
integer $L\in\mathbb{N}$ such that
\begin{equation}
\text{all integers }i\geq L\text{ satisfy }g_{i}\overset{x^{n}}{\equiv}g.
\label{pf.prop.fps.lim.sum-quot.L}%
\end{equation}
Consider these $K$ and $L$. Let us furthermore set $P:=\max\left\{
K,L\right\}  $. Thus, $P\in\mathbb{N}$.

We shall now show that each integer $i\geq P$ satisfies $\left[  x^{n}\right]
\dfrac{f_{i}}{g_{i}}=\left[  x^{n}\right]  \dfrac{f}{g}$.

Indeed, let $i\geq P$ be an integer. Then, (\ref{pf.prop.fps.lim.sum-quot.K})
yields $f_{i}\overset{x^{n}}{\equiv}f$ (since $i\geq P=\max\left\{
K,L\right\}  \geq K$), whereas (\ref{pf.prop.fps.lim.sum-quot.L}) yields
$g_{i}\overset{x^{n}}{\equiv}g$ (since $i\geq P=\max\left\{  K,L\right\}  \geq
L$). Since both FPSs $g_{i}$ and $g$ are invertible (by Claim 1), we thus
obtain
\[
\dfrac{f_{i}}{g_{i}}\overset{x^{n}}{\equiv}\dfrac{f}{g}%
\]
(by Theorem \ref{thm.fps.xneq.props} \textbf{(e)}, applied to $a=f_{i}$,
$b=f$, $c=g_{i}$ and $d=g$). In other words,
\[
\text{each }m\in\left\{  0,1,\ldots,n\right\}  \text{ satisfies }\left[
x^{m}\right]  \dfrac{f_{i}}{g_{i}}=\left[  x^{m}\right]  \dfrac{f}{g}%
\]
(by the definition of $x^{n}$-equivalence). Applying this to $m=n$, we find%
\[
\left[  x^{n}\right]  \dfrac{f_{i}}{g_{i}}=\left[  x^{n}\right]  \dfrac{f}%
{g}.
\]


Forget that we fixed $i$. We thus have shown that all integers $i\geq P$
satisfy $\left[  x^{n}\right]  \dfrac{f_{i}}{g_{i}}=\left[  x^{n}\right]
\dfrac{f}{g}$. Hence, there exists some $N\in\mathbb{N}$ such that%
\[
\text{all integers }i\geq N\text{ satisfy }\left[  x^{n}\right]  \dfrac{f_{i}%
}{g_{i}}=\left[  x^{n}\right]  \dfrac{f}{g}%
\]
(namely, $N=P$). In other words,%
\[
\text{the sequence }\left(  \left[  x^{n}\right]  \dfrac{f_{i}}{g_{i}}\right)
_{i\in\mathbb{N}}\text{ stabilizes to }\left[  x^{n}\right]  \dfrac{f}{g}%
\]
(by the definition of \textquotedblleft stabilizes\textquotedblright).

Forget that we fixed $n$. We thus have shown that for each $n\in\mathbb{N}$,
the sequence $\left(  \left[  x^{n}\right]  \dfrac{f_{i}}{g_{i}}\right)
_{i\in\mathbb{N}}$ stabilizes to $\left[  x^{n}\right]  \dfrac{f}{g}$. In
other words, the sequence $\left(  \dfrac{f_{i}}{g_{i}}\right)  _{i\in
\mathbb{N}}$ coefficientwise stabilizes to $\dfrac{f}{g}$ (by the definition
of \textquotedblleft coefficientwise stabilizing\textquotedblright). In other
words, $\lim\limits_{i\rightarrow\infty}\dfrac{f_{i}}{g_{i}}=\dfrac{f}{g}$.
This proves Proposition \ref{prop.fps.lim.sum-quot} (since we already have
shown that $g$ is invertible).
\end{proof}
\end{fineprint}

\begin{fineprint}
\begin{proof}
[Proof of Theorem \ref{thm.fps.lim.sum-lim}.]The family $\left(  f_{n}\right)
_{n\in\mathbb{N}}$ is summable. In other words, the family $\left(
f_{k}\right)  _{k\in\mathbb{N}}$ is summable (since this is the same family as
$\left(  f_{n}\right)  _{n\in\mathbb{N}}$). In other words, for each
$n\in\mathbb{N}$,
\begin{equation}
\text{all but finitely many }k\in\mathbb{N}\text{ satisfy }\left[
x^{n}\right]  f_{k}=0 \label{pf.thm.fps.lim.sum-lim.1}%
\end{equation}
(by the definition of \textquotedblleft summable\textquotedblright).

Let us define%
\begin{equation}
g_{i}:=\sum_{k=0}^{i}f_{k}\ \ \ \ \ \ \ \ \ \ \text{for each }i\in\mathbb{N}.
\label{pf.thm.fps.lim.sum-lim.gi=}%
\end{equation}
Let us furthermore set%
\begin{equation}
g:=\sum_{k\in\mathbb{N}}f_{k}. \label{pf.thm.fps.lim.sum-lim.g=}%
\end{equation}


Now, fix $n\in\mathbb{N}$. Then, all but finitely many $k\in\mathbb{N}$
satisfy $\left[  x^{n}\right]  f_{k}=0$ (by (\ref{pf.thm.fps.lim.sum-lim.1})).
In other words, there exists a finite subset $J$ of $\mathbb{N}$ such that
\begin{equation}
\text{all }k\in\mathbb{N}\setminus J\text{ satisfy }\left[  x^{n}\right]
f_{k}=0. \label{pf.thm.fps.lim.sum-lim.2}%
\end{equation}
Consider this subset $J$. The set $J$ is a finite set of nonnegative integers,
and thus has an upper bound (since any finite set of nonnegative integers has
an upper bound). In other words, there exists some $m\in\mathbb{N}$ such that%
\begin{equation}
\text{all }k\in J\text{ satisfy }k\leq m. \label{pf.thm.fps.lim.sum-lim.3}%
\end{equation}
Consider this $m$.

Let $k\in\mathbb{N}$ be such that $k\geq m+1$. If we had $k\in J$, then we
would have $k\leq m$ (by (\ref{pf.thm.fps.lim.sum-lim.3})), which would
contradict $k\geq m+1>m$. Thus, we cannot have $k\in J$. Hence, $k\in
\mathbb{N}\setminus J$ (since $k\in\mathbb{N}$ but not $k\in J$). Therefore,
(\ref{pf.thm.fps.lim.sum-lim.2}) yields $\left[  x^{n}\right]  f_{k}=0$.

Forget that we fixed $k$. We thus have shown that%
\begin{equation}
\left[  x^{n}\right]  f_{k}=0\ \ \ \ \ \ \ \ \ \ \text{for each }%
k\in\mathbb{N}\text{ satisfying }k\geq m+1. \label{pf.thm.fps.lim.sum-lim.4}%
\end{equation}


Now, let $i$ be an integer such that $i\geq m$. Then, $i\in\mathbb{N}$ (since
$i\geq m\geq0$) and $i+1\geq m+1$ (since $i\geq m$). From $g=\sum
_{k\in\mathbb{N}}f_{k}$, we obtain%
\begin{align}
\left[  x^{n}\right]  g  &  =\left[  x^{n}\right]  \left(  \sum_{k\in
\mathbb{N}}f_{k}\right)  =\sum_{k\in\mathbb{N}}\left[  x^{n}\right]
f_{k}\ \ \ \ \ \ \ \ \ \ \left(  \text{by (\ref{eq.def.fps.summable.sum}%
)}\right) \nonumber\\
&  =\sum_{k=0}^{i}\left[  x^{n}\right]  f_{k}+\sum_{k=i+1}^{\infty
}\underbrace{\left[  x^{n}\right]  f_{k}}_{\substack{=0\\\text{(by
(\ref{pf.thm.fps.lim.sum-lim.4})}\\\text{(since }k\geq i+1\geq m+1\text{))}%
}}=\sum_{k=0}^{i}\left[  x^{n}\right]  f_{k}+\underbrace{\sum_{k=i+1}^{\infty
}0}_{=0}\nonumber\\
&  =\sum_{k=0}^{i}\left[  x^{n}\right]  f_{k}.
\label{pf.thm.fps.lim.sum-lim.5}%
\end{align}
On the other hand, from $g_{i}=\sum_{k=0}^{i}f_{k}$, we obtain%
\[
\left[  x^{n}\right]  g_{i}=\left[  x^{n}\right]  \left(  \sum_{k=0}^{i}%
f_{k}\right)  =\sum_{k=0}^{i}\left[  x^{n}\right]  f_{k}.
\]
Comparing this with (\ref{pf.thm.fps.lim.sum-lim.5}), we obtain $\left[
x^{n}\right]  g_{i}=\left[  x^{n}\right]  g$.

Forget that we fixed $i$. We thus have shown that
\[
\text{all integers }i\geq m\text{ satisfy }\left[  x^{n}\right]  g_{i}=\left[
x^{n}\right]  g.
\]
Hence, there exists some $N\in\mathbb{N}$ such that
\[
\text{all integers }i\geq N\text{ satisfy }\left[  x^{n}\right]  g_{i}=\left[
x^{n}\right]  g
\]
(namely, $N=m$). In other words, the sequence $\left(  \left[  x^{n}\right]
g_{i}\right)  _{i\in\mathbb{N}}$ stabilizes to $\left[  x^{n}\right]  g$ (by
the definition of \textquotedblleft stabilizes\textquotedblright).

Forget that we fixed $n$. We thus have shown that for each $n\in\mathbb{N}$,%
\[
\text{the sequence }\left(  \left[  x^{n}\right]  g_{i}\right)  _{i\in
\mathbb{N}}\text{ stabilizes to }\left[  x^{n}\right]  g.
\]
In other words, the sequence $\left(  g_{i}\right)  _{i\in\mathbb{N}}$
coefficientwise stabilizes to $g$. In other words, $\lim\limits_{i\rightarrow
\infty}g_{i}=g$. In view of (\ref{pf.thm.fps.lim.sum-lim.gi=}) and
(\ref{pf.thm.fps.lim.sum-lim.g=}), we can rewrite this as%
\[
\lim\limits_{i\rightarrow\infty}\sum_{k=0}^{i}f_{k}=\sum_{k\in\mathbb{N}}%
f_{k}.
\]
Renaming the summation index $k$ as $n$ everywhere in this equality, we can
rewrite it as%
\[
\lim\limits_{i\rightarrow\infty}\sum_{n=0}^{i}f_{n}=\sum_{n\in\mathbb{N}}%
f_{n}.
\]
This proves Theorem \ref{thm.fps.lim.sum-lim}.
\end{proof}

\begin{proof}
[Proof of Theorem \ref{thm.fps.lim.prod-lim}.]The family $\left(
f_{n}\right)  _{n\in\mathbb{N}}$ is multipliable. In other words, the family
$\left(  f_{k}\right)  _{k\in\mathbb{N}}$ is multipliable (since this is the
same family as $\left(  f_{n}\right)  _{n\in\mathbb{N}}$). In other words,
each coefficient in the product of this family is finitely determined (by the
definition of \textquotedblleft multipliable\textquotedblright).

Let us define%
\begin{equation}
g_{i}:=\prod_{k=0}^{i}f_{k}\ \ \ \ \ \ \ \ \ \ \text{for each }i\in\mathbb{N}.
\label{pf.thm.fps.lim.prod-lim.gi=}%
\end{equation}
Let us furthermore set%
\begin{equation}
g:=\prod_{k\in\mathbb{N}}f_{k}. \label{pf.thm.fps.lim.prod-lim.g=}%
\end{equation}


Now, fix $n\in\mathbb{N}$. Then, the $x^{n}$-coefficient in the product of the
family $\left(  f_{k}\right)  _{k\in\mathbb{N}}$ is finitely determined (since
each coefficient in the product of this family is finitely determined). In
other words, there is a finite subset $M$ of $\mathbb{N}$ that determines the
$x^{n}$-coefficient in the product of $\left(  f_{k}\right)  _{k\in\mathbb{N}%
}$ (by the definition of \textquotedblleft finitely
determined\textquotedblright). Consider this subset $M$.

The set $M$ is a finite set of nonnegative integers, and thus has an upper
bound (since any finite set of nonnegative integers has an upper bound). In
other words, there exists some $m\in\mathbb{N}$ such that%
\begin{equation}
\text{all }k\in M\text{ satisfy }k\leq m. \label{pf.thm.fps.lim.prod-lim.3}%
\end{equation}
Consider this $m$.

Now, let $i$ be an integer be such that $i\geq m$. Then, $i\in\mathbb{N}$
(since $i\geq m\geq0$) and $m\leq i$. From (\ref{pf.thm.fps.lim.prod-lim.3}),
we see that all $k\in M$ satisfy $k\leq m\leq i$ and therefore $k\in\left\{
0,1,\ldots,i\right\}  $. In other words, $M\subseteq\left\{  0,1,\ldots
,i\right\}  $.

However, the set $M$ determines the $x^{n}$-coefficient in the product of
$\left(  f_{k}\right)  _{k\in\mathbb{N}}$. In other words, every finite subset
$J$ of $\mathbb{N}$ satisfying $M\subseteq J\subseteq\mathbb{N}$ satisfies%
\[
\left[  x^{n}\right]  \left(  \prod_{k\in J}f_{k}\right)  =\left[
x^{n}\right]  \left(  \prod_{k\in M}f_{k}\right)
\]
(by the definition of \textquotedblleft determines the $x^{n}$-coefficient in
the product of $\left(  f_{k}\right)  _{k\in\mathbb{N}}$\textquotedblright).
Applying this to $J=\left\{  0,1,\ldots,i\right\}  $, we obtain%
\begin{equation}
\left[  x^{n}\right]  \left(  \prod_{k\in\left\{  0,1,\ldots,i\right\}  }%
f_{k}\right)  =\left[  x^{n}\right]  \left(  \prod_{k\in M}f_{k}\right)
\label{pf.thm.fps.lim.prod-lim.4}%
\end{equation}
(since $\left\{  0,1,\ldots,i\right\}  $ is a finite subset of $\mathbb{N}$
satisfying $M\subseteq\left\{  0,1,\ldots,i\right\}  \subseteq\mathbb{N}$).

Moreover, the definition of the product of the multipliable family $\left(
f_{k}\right)  _{k\in\mathbb{N}}$ (Definition \ref{def.fps.multipliable}
\textbf{(b)}) shows that%
\[
\left[  x^{n}\right]  \left(  \prod_{k\in\mathbb{N}}f_{k}\right)  =\left[
x^{n}\right]  \left(  \prod_{k\in M}f_{k}\right)
\]
(since $M$ is a finite subset of $\mathbb{N}$ that determines the $x^{n}%
$-coefficient in the product of $\left(  f_{k}\right)  _{k\in\mathbb{N}}$).
Comparing this with (\ref{pf.thm.fps.lim.prod-lim.4}), we find%
\[
\left[  x^{n}\right]  \left(  \prod_{k\in\left\{  0,1,\ldots,i\right\}  }%
f_{k}\right)  =\left[  x^{n}\right]  \left(  \prod_{k\in\mathbb{N}}%
f_{k}\right)  .
\]


In view of $\prod_{k\in\left\{  0,1,\ldots,i\right\}  }f_{k}=\prod_{k=0}%
^{i}f_{k}=g_{i}$ (by (\ref{pf.thm.fps.lim.prod-lim.gi=})) and $\prod
_{k\in\mathbb{N}}f_{k}=g$ (by (\ref{pf.thm.fps.lim.prod-lim.g=})), we can
rewrite this as%
\[
\left[  x^{n}\right]  g_{i}=\left[  x^{n}\right]  g.
\]


Forget that we fixed $i$. We thus have shown that
\[
\text{all integers }i\geq m\text{ satisfy }\left[  x^{n}\right]  g_{i}=\left[
x^{n}\right]  g.
\]
Hence, there exists some $N\in\mathbb{N}$ such that
\[
\text{all integers }i\geq N\text{ satisfy }\left[  x^{n}\right]  g_{i}=\left[
x^{n}\right]  g
\]
(namely, $N=m$). In other words, the sequence $\left(  \left[  x^{n}\right]
g_{i}\right)  _{i\in\mathbb{N}}$ stabilizes to $\left[  x^{n}\right]  g$ (by
the definition of \textquotedblleft stabilizes\textquotedblright).

Forget that we fixed $n$. We thus have shown that for each $n\in\mathbb{N}$,%
\[
\text{the sequence }\left(  \left[  x^{n}\right]  g_{i}\right)  _{i\in
\mathbb{N}}\text{ stabilizes to }\left[  x^{n}\right]  g.
\]
In other words, the sequence $\left(  g_{i}\right)  _{i\in\mathbb{N}}$
coefficientwise stabilizes to $g$. In other words, $\lim\limits_{i\rightarrow
\infty}g_{i}=g$. In view of (\ref{pf.thm.fps.lim.prod-lim.g=}) and
(\ref{pf.thm.fps.lim.prod-lim.g=}), we can rewrite this as%
\[
\lim\limits_{i\rightarrow\infty}\prod_{k=0}^{i}f_{k}=\prod_{k\in\mathbb{N}%
}f_{k}.
\]
Renaming the product index $k$ as $n$ everywhere in this equality, we can
rewrite it as%
\[
\lim\limits_{i\rightarrow\infty}\prod_{n=0}^{i}f_{n}=\prod_{n\in\mathbb{N}%
}f_{n}.
\]
This proves Theorem \ref{thm.fps.lim.prod-lim}.
\end{proof}

\begin{proof}
[Proof of Corollary \ref{cor.fps.lim.fps-as-pol}.]Let a FPS $a\in K\left[
\left[  x\right]  \right]  $ be written in the form $a=\sum_{n\in\mathbb{N}%
}a_{n}x^{n}$ with $a_{n}\in K$. Then, Corollary \ref{cor.fps.sumakxk} shows
that the family $\left(  a_{n}x^{n}\right)  _{n\in\mathbb{N}}$ is summable.
Hence, Theorem \ref{thm.fps.lim.sum-lim} (applied to $f_{n}=a_{n}x^{n}$)
yields
\[
\lim\limits_{i\rightarrow\infty}\sum_{n=0}^{i}a_{n}x^{n}=\sum_{n\in\mathbb{N}%
}a_{n}x^{n}=a.
\]
In other words, $a=\lim\limits_{i\rightarrow\infty}\sum_{n=0}^{i}a_{n}x^{n}$.
This proves Corollary \ref{cor.fps.lim.fps-as-pol}.
\end{proof}
\end{fineprint}

\begin{fineprint}
\begin{proof}
[Proof of Theorem \ref{thm.fps.lim.sum-lim-conv}.]Let us define%
\begin{equation}
g_{i}:=\sum_{k=0}^{i}f_{k}\ \ \ \ \ \ \ \ \ \ \text{for each }i\in\mathbb{N}.
\label{pf.thm.fps.lim.sum-lim-conv.gi=}%
\end{equation}
We have assumed that the limit $\lim\limits_{i\rightarrow\infty}\sum_{n=0}%
^{i}f_{n}$ exists. Let us denote this limit by $g$. Thus,%
\begin{align}
g  &  =\lim\limits_{i\rightarrow\infty}\sum_{n=0}^{i}f_{n}=\lim
\limits_{i\rightarrow\infty}\underbrace{\sum_{k=0}^{i}f_{k}}_{\substack{=g_{i}%
\\\text{(by (\ref{pf.thm.fps.lim.sum-lim-conv.gi=}))}}%
}\ \ \ \ \ \ \ \ \ \ \left(
\begin{array}
[c]{c}%
\text{here, we have renamed the}\\
\text{summation index }n\text{ as }k
\end{array}
\right) \nonumber\\
&  =\lim\limits_{i\rightarrow\infty}g_{i}.
\label{pf.thm.fps.lim.sum-lim-conv.g=lim}%
\end{align}
In other words, the sequence $\left(  g_{i}\right)  _{i\in\mathbb{N}}$
coefficientwise stabilizes to $g$. In other words, for each $n\in\mathbb{N}$,%
\begin{equation}
\text{the sequence }\left(  \left[  x^{n}\right]  g_{i}\right)  _{i\in
\mathbb{N}}\text{ stabilizes to }\left[  x^{n}\right]  g.
\label{pf.thm.fps.lim.sum-lim-conv.gstab}%
\end{equation}


Let $n\in\mathbb{N}$. Then, the sequence $\left(  \left[  x^{n}\right]
g_{i}\right)  _{i\in\mathbb{N}}$ stabilizes to $\left[  x^{n}\right]  g$ (by
(\ref{pf.thm.fps.lim.sum-lim-conv.gstab})). In other words, there exists some
$N\in\mathbb{N}$ such that
\begin{equation}
\text{all integers }i\geq N\text{ satisfy }\left[  x^{n}\right]  g_{i}=\left[
x^{n}\right]  g. \label{pf.thm.fps.lim.sum-lim-conv.1}%
\end{equation}
Consider this $N$. Let $M$ be the set $\left\{  0,1,\ldots,N\right\}  $; this
is a finite subset of $\mathbb{N}$.

Let $i\in\mathbb{N}\setminus M$. Thus,
\begin{align*}
i  &  \in\mathbb{N}\setminus M=\mathbb{N}\setminus\left\{  0,1,\ldots
,N\right\}  \ \ \ \ \ \ \ \ \ \ \left(  \text{since }M=\left\{  0,1,\ldots
,N\right\}  \right) \\
&  =\left\{  N+1,\ N+2,\ N+3,\ \ldots\right\}  .
\end{align*}
In other words, $i$ is a nonnegative integer with $i\geq N+1$. Hence, $i-1\geq
N$. Therefore, (\ref{pf.thm.fps.lim.sum-lim-conv.1}) (applied to $i-1$ instead
of $i$) yields $\left[  x^{n}\right]  g_{i-1}=\left[  x^{n}\right]  g$. But
(\ref{pf.thm.fps.lim.sum-lim-conv.1}) also yields $\left[  x^{n}\right]
g_{i}=\left[  x^{n}\right]  g$ (since $i\geq N+1\geq N$). Comparing these two
equalities, we find
\begin{equation}
\left[  x^{n}\right]  g_{i-1}=\left[  x^{n}\right]  g_{i}.
\label{pf.thm.fps.lim.sum-lim-conv.2}%
\end{equation}


However, (\ref{pf.thm.fps.lim.sum-lim-conv.gi=}) yields
\begin{equation}
g_{i}=\sum_{k=0}^{i}f_{k}=f_{i}+\sum_{k=0}^{i-1}f_{k}.
\label{pf.thm.fps.lim.sum-lim-conv.3}%
\end{equation}
Moreover, (\ref{pf.thm.fps.lim.sum-lim-conv.gi=}) (applied to $i-1$ instead of
$i$) yields $g_{i-1}=\sum_{k=0}^{i-1}f_{k}$. In view of this, we can rewrite
(\ref{pf.thm.fps.lim.sum-lim-conv.3}) as%
\[
g_{i}=f_{i}+g_{i-1}.
\]
Hence, $\left[  x^{n}\right]  g_{i}=\left[  x^{n}\right]  \left(
f_{i}+g_{i-1}\right)  =\left[  x^{n}\right]  f_{i}+\underbrace{\left[
x^{n}\right]  g_{i-1}}_{\substack{=\left[  x^{n}\right]  g_{i}\\\text{(by
(\ref{pf.thm.fps.lim.sum-lim-conv.2}))}}}=\left[  x^{n}\right]  f_{i}+\left[
x^{n}\right]  g_{i}$. Subtracting $\left[  x^{n}\right]  g_{i}$ from both
sides of this equality, we find $0=\left[  x^{n}\right]  f_{i}$. In other
words, $\left[  x^{n}\right]  f_{i}=0$.

Forget that we fixed $i$. We thus have shown that all $i\in\mathbb{N}\setminus
M$ satisfy $\left[  x^{n}\right]  f_{i}=0$. Hence, all but finitely many
$i\in\mathbb{N}$ satisfy $\left[  x^{n}\right]  f_{i}=0$ (since all but
finitely many $i\in\mathbb{N}$ belong to $\mathbb{N}\setminus M$ (because the
set $M$ is finite)). Renaming the index $i$ as $k$ in this statement, we
obtain the following: All but finitely many $k\in\mathbb{N}$ satisfy $\left[
x^{n}\right]  f_{k}=0$.

Forget that we fixed $n$. We thus have shown that for each $n\in\mathbb{N}$,
\[
\text{all but finitely many }k\in\mathbb{N}\text{ satisfy }\left[
x^{n}\right]  f_{k}=0.
\]
In other words, the family $\left(  f_{n}\right)  _{n\in\mathbb{N}}$ is
summable (by the definition of \textquotedblleft summable\textquotedblright).

Hence, Theorem \ref{thm.fps.lim.sum-lim} yields $\lim\limits_{i\rightarrow
\infty}\sum_{n=0}^{i}f_{n}=\sum_{n\in\mathbb{N}}f_{n}$. In other words,
$\sum_{n\in\mathbb{N}}f_{n}=\lim\limits_{i\rightarrow\infty}\sum_{n=0}%
^{i}f_{n}$. This completes the proof of Theorem \ref{thm.fps.lim.sum-lim-conv}.
\end{proof}

\begin{proof}
[Proof of Theorem \ref{thm.fps.lim.prod-lim-conv}.]Let us define%
\begin{equation}
g_{i}:=\prod_{k=0}^{i}f_{k}\ \ \ \ \ \ \ \ \ \ \text{for each }i\in\mathbb{N}.
\label{pf.thm.fps.lim.prod-lim-conv.gi=}%
\end{equation}


We have assumed that the limit $\lim\limits_{i\rightarrow\infty}\prod
_{n=0}^{i}f_{n}$ exists. Let us denote this limit by $g$. Thus,%
\begin{align}
g  &  =\lim\limits_{i\rightarrow\infty}\prod_{n=0}^{i}f_{n}=\lim
\limits_{i\rightarrow\infty}\underbrace{\prod_{k=0}^{i}f_{k}}%
_{\substack{=g_{i}\\\text{(by (\ref{pf.thm.fps.lim.prod-lim-conv.gi=}))}%
}}\ \ \ \ \ \ \ \ \ \ \left(
\begin{array}
[c]{c}%
\text{here, we have renamed the}\\
\text{product index }n\text{ as }k
\end{array}
\right) \nonumber\\
&  =\lim\limits_{i\rightarrow\infty}g_{i}.
\label{pf.thm.fps.lim.prod-lim-conv.g=lim}%
\end{align}


Let $n\in\mathbb{N}$ be arbitrary. Lemma \ref{lem.fps.lim.xn-equiv} (applied
to $g_{i}$ and $g$ instead of $f_{i}$ and $f$) shows that there exists some
integer $N\in\mathbb{N}$ such that
\begin{equation}
\text{all integers }i\geq N\text{ satisfy }g_{i}\overset{x^{n}}{\equiv}g
\label{pf.thm.fps.lim.prod-lim-conv.gi-equiv-g}%
\end{equation}
(since $g=\lim\limits_{i\rightarrow\infty}g_{i}$). Consider this $N$.

Let $M=\left\{  0,1,\ldots,N\right\}  $. Thus, $M$ is a finite subset of
$\mathbb{N}$. We shall show that this subset $M$ determines the $x^{n}%
$-coefficient in the product of $\left(  f_{k}\right)  _{k\in\mathbb{N}}$.

For this purpose, we first observe that $M=\left\{  0,1,\ldots,N\right\}  $,
so that%
\begin{equation}
\prod_{k\in M}f_{k}=\prod_{k\in\left\{  0,1,\ldots,N\right\}  }f_{k}%
=\prod_{k=0}^{N}f_{k}=g_{N} \label{pf.thm.fps.lim.prod-lim-conv.=gN}%
\end{equation}
(since $g_{N}$ was defined to be $\prod_{k=0}^{N}f_{k}$). Applying
(\ref{pf.thm.fps.lim.prod-lim-conv.gi-equiv-g}) to $i=N$, we obtain
$g_{N}\overset{x^{n}}{\equiv}g$ (since $N\geq N$). Therefore, $g\overset{x^{n}%
}{\equiv}g_{N}$ (since $\overset{x^{n}}{\equiv}$ is an equivalence relation
and thus symmetric).

Next, we shall show two claims:

\begin{statement}
\textit{Claim 1:} For each integer $j>N$, we have $g\overset{x^{n}}{\equiv
}gf_{j}$.
\end{statement}

\begin{proof}
[Proof of Claim 1.]Let $j>N$ be an integer. Thus, $j\geq N+1$, so that
$j-1\geq N$. Hence, (\ref{pf.thm.fps.lim.prod-lim-conv.gi-equiv-g}) (applied
to $i=j-1$) yields $g_{j-1}\overset{x^{n}}{\equiv}g$. However,
(\ref{pf.thm.fps.lim.prod-lim-conv.gi-equiv-g}) (applied to $i=j$) yields
$g_{j}\overset{x^{n}}{\equiv}g$ (since $j\geq N+1\geq N$). Thus,
$g\overset{x^{n}}{\equiv}g_{j}$ (since $\overset{x^{n}}{\equiv}$ is an
equivalence relation and thus symmetric).

Also, we obviously have $f_{j}\overset{x^{n}}{\equiv}f_{j}$ (since
$\overset{x^{n}}{\equiv}$ is an equivalence relation and thus reflexive).

However, the definition of $g_{j-1}$ yields $g_{j-1}=\prod_{k=0}^{j-1}f_{k}$.
Furthermore, the definition of $g_{j}$ yields%
\[
g_{j}=\prod_{k=0}^{j}f_{k}=\underbrace{\left(  \prod_{k=0}^{j-1}f_{k}\right)
}_{=g_{j-1}}f_{j}=g_{j-1}f_{j}.
\]
But recall that $g_{j-1}\overset{x^{n}}{\equiv}g$ and $f_{j}\overset{x^{n}%
}{\equiv}f_{j}$. Hence, $g_{j-1}f_{j}\overset{x^{n}}{\equiv}gf_{j}$ (by
(\ref{eq.thm.fps.xneq.props.b.*}), applied to $a=g_{j-1}$, $b=g$, $c=f_{j}$
and $d=f_{j}$). In view of $g_{j}=g_{j-1}f_{j}$, we can rewrite this as
$g_{j}\overset{x^{n}}{\equiv}gf_{j}$.

Now, recall that $g\overset{x^{n}}{\equiv}g_{j}$. Thus, $g\overset{x^{n}%
}{\equiv}g_{j}\overset{x^{n}}{\equiv}gf_{j}$, so that $g\overset{x^{n}%
}{\equiv}gf_{j}$ (since $\overset{x^{n}}{\equiv}$ is an equivalence relation
and thus transitive). This proves Claim 1.
\end{proof}

\begin{statement}
\textit{Claim 2:} If $U$ is any finite subset of $\mathbb{N}$ satisfying
$M\subseteq U$, then $g\overset{x^{n}}{\equiv}\prod_{k\in U}f_{k}$.
\end{statement}

\begin{proof}
[Proof of Claim 2.]We induct on the nonnegative integer $\left\vert U\setminus
M\right\vert $:

\textit{Base case:} Let us check that Claim 2 holds when $\left\vert
U\setminus M\right\vert =0$.

Indeed, let $U$ be any finite subset of $\mathbb{N}$ satisfying $M\subseteq U$
and $\left\vert U\setminus M\right\vert =0$. Then, $U\setminus M=\varnothing$
(since $\left\vert U\setminus M\right\vert =0$), so that $U\subseteq M$.
Combined with $M\subseteq U$, this yields $U=M$. Thus, $\prod_{k\in U}%
f_{k}=\prod_{k\in M}f_{k}=g_{N}$ (by (\ref{pf.thm.fps.lim.prod-lim-conv.=gN}%
)). But we have already proved that $g\overset{x^{n}}{\equiv}g_{N}$. In other
words, $g\overset{x^{n}}{\equiv}\prod_{k\in U}f_{k}$ (since $\prod_{k\in
U}f_{k}=g_{N}$). Thus, Claim 2 is proved under the assumption that $\left\vert
U\setminus M\right\vert =0$. This completes the base case.

\textit{Induction step:} Let $s\in\mathbb{N}$. Assume (as the induction
hypothesis) that Claim 2 holds when $\left\vert U\setminus M\right\vert =s$.
We must now prove that Claim 2 holds when $\left\vert U\setminus M\right\vert
=s+1$.

Indeed, let $U$ be any finite subset of $\mathbb{N}$ satisfying $M\subseteq U$
and $\left\vert U\setminus M\right\vert =s+1$. Then, the set $U\setminus M$ is
nonempty (since $\left\vert U\setminus M\right\vert =s+1>s\geq0$). Hence,
there exists some $u\in U\setminus M$. Consider this $u$. From $u\in
U\setminus M$, we obtain $\left\vert \left(  U\setminus M\right)
\setminus\left\{  u\right\}  \right\vert =\left\vert U\setminus M\right\vert
-1=s$ (since $\left\vert U\setminus M\right\vert =s+1$). In view of $\left(
U\setminus M\right)  \setminus\left\{  u\right\}  =\left(  U\setminus\left\{
u\right\}  \right)  \setminus M$, we can rewrite this as $\left\vert \left(
U\setminus\left\{  u\right\}  \right)  \setminus M\right\vert =s$. Moreover,
$U\setminus\left\{  u\right\}  $ is a finite subset of $\mathbb{N}$ (since $U$
is a finite subset of $\mathbb{N}$), and satisfies $M\subseteq U\setminus
\left\{  u\right\}  $ (since $u\in U\setminus M$ and thus $u\notin M$ and
therefore $M\setminus\left\{  u\right\}  =M$, so that $M=\underbrace{M}%
_{\subseteq U}\setminus\left\{  u\right\}  \subseteq U\setminus\left\{
u\right\}  $). Hence, by the induction hypothesis, we can apply Claim 2 to
$U\setminus\left\{  u\right\}  $ instead of $U$ (since $\left\vert \left(
U\setminus\left\{  u\right\}  \right)  \setminus M\right\vert =s$). As a
result, we obtain $g\overset{x^{n}}{\equiv}\prod_{k\in U\setminus\left\{
u\right\}  }f_{k}$.

However, $u\in U\setminus M\subseteq U$, so that $\prod_{k\in U}f_{k}=\left(
\prod_{k\in U\setminus\left\{  u\right\}  }f_{k}\right)  f_{u}$ (here, we have
split off the factor for $k=u$ from the product).

We have $g\overset{x^{n}}{\equiv}\prod_{k\in U\setminus\left\{  u\right\}
}f_{k}$ (as we have proved above) and $f_{u}\overset{x^{n}}{\equiv}f_{u}$
(since $\overset{x^{n}}{\equiv}$ is an equivalence relation and thus
reflexive). Hence, $gf_{u}\overset{x^{n}}{\equiv}\left(  \prod_{k\in
U\setminus\left\{  u\right\}  }f_{k}\right)  f_{u}$ (by
(\ref{eq.thm.fps.xneq.props.b.*}), applied to $a=g$, $b=\prod_{k\in
U\setminus\left\{  u\right\}  }f_{k}$, $c=f_{u}$ and $d=f_{u}$). In view of
$\prod_{k\in U}f_{k}=\left(  \prod_{k\in U\setminus\left\{  u\right\}  }%
f_{k}\right)  f_{u}$, we can rewrite this as $gf_{u}\overset{x^{n}}{\equiv
}\prod_{k\in U}f_{k}$.

But $u\in\underbrace{U}_{\subseteq\mathbb{N}}\setminus\underbrace{M}%
_{=\left\{  0,1,\ldots,N\right\}  }\subseteq\mathbb{N}\setminus\left\{
0,1,\ldots,N\right\}  =\left\{  N+1,\ N+2,\ N+3,\ \ldots\right\}  $, so that
$u\geq N+1>N$. Hence, Claim 1 (applied to $j=u$) yields $g\overset{x^{n}%
}{\equiv}gf_{u}$. Therefore, $g\overset{x^{n}}{\equiv}gf_{u}\overset{x^{n}%
}{\equiv}\prod_{k\in U}f_{k}$. Therefore, $g\overset{x^{n}}{\equiv}\prod_{k\in
U}f_{k}$ (since $\overset{x^{n}}{\equiv}$ is an equivalence relation and thus transitive).

Forget that we fixed $U$. We thus have shown that if $U$ is any finite subset
of $\mathbb{N}$ satisfying $M\subseteq U$ and $\left\vert U\setminus
M\right\vert =s+1$, then $g\overset{x^{n}}{\equiv}\prod_{k\in U}f_{k}$. In
other words, Claim 2 holds when $\left\vert U\setminus M\right\vert =s+1$.
This completes the induction step. Thus, Claim 2 is proved by induction.
\end{proof}

Now, let $J$ be a finite subset of $\mathbb{N}$ satisfying $M\subseteq
J\subseteq\mathbb{N}$. Then, $g\overset{x^{n}}{\equiv}\prod_{k\in J}f_{k}$ (by
Claim 2, applied to $U=J$). In other words,%
\[
\text{each }m\in\left\{  0,1,\ldots,n\right\}  \text{ satisfies }\left[
x^{m}\right]  g=\left[  x^{m}\right]  \left(  \prod_{k\in J}f_{k}\right)
\]
(by the definition of $x^{n}$-equivalence). Applying this to $m=n$, we find%
\[
\left[  x^{n}\right]  g=\left[  x^{n}\right]  \left(  \prod_{k\in J}%
f_{k}\right)  .
\]
The same argument can be applied to $M$ instead of $J$ (since $M\subseteq
M\subseteq\mathbb{N}$), and thus yields%
\[
\left[  x^{n}\right]  g=\left[  x^{n}\right]  \left(  \prod_{k\in M}%
f_{k}\right)  .
\]
Comparing these two equalities, we find
\[
\left[  x^{n}\right]  \left(  \prod_{k\in J}f_{k}\right)  =\left[
x^{n}\right]  \left(  \prod_{k\in M}f_{k}\right)  .
\]


Forget that we fixed $J$. We have now shown that every finite subset $J$ of
$\mathbb{N}$ satisfying $M\subseteq J\subseteq\mathbb{N}$ satisfies%
\[
\left[  x^{n}\right]  \left(  \prod_{k\in J}f_{k}\right)  =\left[
x^{n}\right]  \left(  \prod_{k\in M}f_{k}\right)  .
\]
In other words, the set $M$ determines the $x^{n}$-coefficient in the product
of $\left(  f_{k}\right)  _{k\in\mathbb{N}}$ (by the definition of
\textquotedblleft determines the $x^{n}$-coefficient in the product of
$\left(  f_{k}\right)  _{k\in\mathbb{N}}$\textquotedblright). Hence, the
$x^{n}$-coefficient in the product of $\left(  f_{k}\right)  _{k\in\mathbb{N}%
}$ is finitely determined (by the definition of \textquotedblleft finitely
determined\textquotedblright, since $M$ is a finite subset of $\mathbb{N}$).

Forget that we fixed $n$. We thus have shown that for each $n\in\mathbb{N}$,
the $x^{n}$-coefficient in the product of $\left(  f_{k}\right)
_{k\in\mathbb{N}}$ is finitely determined. In other words, each coefficient in
the product of $\left(  f_{k}\right)  _{k\in\mathbb{N}}$ is finitely
determined. In other words, the family $\left(  f_{k}\right)  _{k\in
\mathbb{N}}$ is multipliable (by the definition of \textquotedblleft
multipliable\textquotedblright). In other words, the family $\left(
f_{n}\right)  _{n\in\mathbb{N}}$ is multipliable (since this is the same
family as $\left(  f_{k}\right)  _{k\in\mathbb{N}}$).

Hence, Theorem \ref{thm.fps.lim.prod-lim} yields $\lim\limits_{i\rightarrow
\infty}\prod_{n=0}^{i}f_{n}=\prod_{n\in\mathbb{N}}f_{n}$. In other words,
$\prod_{n\in\mathbb{N}}f_{n}=\lim\limits_{i\rightarrow\infty}\prod_{n=0}%
^{i}f_{n}$. This completes the proof of Theorem
\ref{thm.fps.lim.prod-lim-conv}.
\end{proof}
\end{fineprint}
